Electrospray and colloid thrusters rely on the formation of a Taylor cone, an electrohydrodynamic interface where electric Maxwell stress balances capillary stress. High-fidelity electrospray modeling can require expensive multiphysics software, dense meshes, and substantial computation, creating a barrier for early-stage design and student research. This project develops a lightweight, open-source, two-dimensional axisymmetric Python solver for onset-level Taylor-cone modeling on consumer hardware. The solver assembles sparse finite-difference Laplace/Poisson electrostatic systems, reconstructs electric fields, evaluates Maxwell and capillary pressure on a candidate interface, computes a Young--Laplace--Maxwell residual, and optionally includes effective Gaussian or threshold-activated space-charge shielding closures. Preliminary numerical validation demonstrates recovery of the classical Taylor half-angle benchmark, manufactured Poisson convergence, voltage scaling consistency, residual sanity checks, and a passing automated test suite. The planned experimental validation will image supervised Taylor-cone formation, extract cone profiles and half-angles using computer vision, and compare measured profiles and onset-voltage trends against solver predictions using angle error, profile RMSE, and voltage-error metrics. The project is intentionally framed as a reduced-order onset model, not a full CFD, cone-jet breakup, or complete thruster-performance simulator. Its goal is to evaluate how accurately a transparent sparse solver can capture key Taylor-cone observables while remaining accessible to students and small laboratories.
